\section{Related Work}
\label{sec:RelatedWork}

\textbf{Reference-based metric methods}~\cite{meteor, bleu, Spice, cider, Faier} evaluate caption quality by comparing generated captions with ground-truth captions. BLEU~\cite{bleu} is a fast, cost-effective, and language-independent metric for machine translation, correlating well with human assessments. METEOR~\cite{meteor} also evaluates machine translation and shows a high correlation with human judgments, significantly outperforming BLEU. CIDEr~\cite{cider} measures the similarity of generated sentences against a set of human-written ground-truth sentences, serving as an automatic consensus metric for image description quality. However, these metrics primarily focus on n-gram overlap, which is neither necessary nor sufficient for simulating human judgment. To address these limitations, SPICE introduces a metric based on semantic propositional content over scene graphs, though its scores highly rely on the format of reference captions.  Newer methods like FAIEr~\cite{Faier} leverage visual and textual scene graphs for a more robust evaluation.

\textbf{Reference-free metric methods}~\cite{infometic, Clipscore, Tiger} utilize semantic vectors from the reference image to assess caption similarity. InfoMetIC~\cite{infometic} is an informative metric for reference-free image caption evaluation, capable of identifying incorrect words and unmentioned image regions with fine-grained precision. It provides a text precision score, a vision recall score, and an overall quality score at a coarse-grained level, with the latter showing significantly better correlation with human judgments than existing metrics across multiple benchmarks. CLIPScore~\cite{Clipscore} employs CLIP~\cite{clip} for robust automatic evaluation of image captioning without references, focusing on image-text compatibility. Although complementary to reference-based metrics that emphasize text-text similarities, CLIPScore is relatively weaker for tasks requiring richer contextual knowledge, such as news captions. TIGEr~\cite{Tiger} assesses caption quality by evaluating both the representation of image content and the alignment of machine-generated captions with human-generated ones.


\textbf{Question-based evaluation methods}~\cite{QACE, cho2023davidsonian, playgroundv3} have been developed to enhance the assessment of caption quality in visual generation tasks. The QACE framework~\cite{QACE} generates questions from captions to evaluate their quality. A similar approach has been proposed for image generation models, where the Davidsonian Scene Graph (DSG)~\cite{cho2023davidsonian} organizes questions into dependency graphs, facilitating comprehensive evaluation of text-to-image models. Inspired by DSG and DPG-bench~\cite{Ella}, Playground v3~\cite{playgroundv3} introduces CapsBench, a benchmark for image captioning that uses ``yes-no" question-answer pairs. However, these benchmarks focus solely on image captioning, while video captioning requires consideration of additional factors such as actions, motion, camera movements, and shot types, which are crucial for text-to-video generation.